% ============================================================
% 片段：s00  题名页、摘要、译者说明、目录
% 对应原文 PDF 第 1--2 页（题名页、目录）
% ============================================================

\maketitle
\thispagestyle{empty}

\begin{center}
  {\small 原文：arXiv:1207.0496v1 [hep-th], 2 Jul 2012, 44 pp.}\\[2pt]
  {\small 2012 年讲于清华大学数学科学中心（Mathematical Sciences Center, Tsinghua University）}\\[2pt]
  {\small 原题：\emph{Lectures on Mirror Symmetry and Topological String Theory}}
\end{center}

\vspace{1.2em}
\newpage

% ---------------------------------------------------------------------------
\section*{摘要}
\addcontentsline{toc}{section}{摘要}
% ---------------------------------------------------------------------------

本讲义是 2012 年在清华大学数学科学中心所做的一系列关于镜像对称（mirror symmetry）与拓扑弦理论
（topological string theory）的讲座笔记。文中回顾了 $\mathcal{N}=2$ 超共形代数、其形变以及手征环
（chiral ring）。可以构造一个拓扑场论，其可观测量仅为手征环的元素；将其与引力耦合，便得到拓扑弦理论。
镜像对称把拓扑弦的 A 模型与 B 模型等同起来，从而在数学中引出令人惊讶的联系，并提供精确计算的工具
以及物理上的新洞见。最后回顾了把拓扑弦理论的高亏格振幅表示为多项式的递推构造。

% ---------------------------------------------------------------------------
\section*{译者说明}
\addcontentsline{toc}{section}{译者说明}
% ---------------------------------------------------------------------------

\begin{enumerate}[leftmargin=1.8em,itemsep=2pt,topsep=4pt]
  \item \textbf{底本}：本译本依据 arXiv:1207.0496v1（Murad Alim，哈佛大学杰斐逊物理实验室；
        2012 年 7 月 2 日提交，共 44 页）全文译出，正文、脚注、图表标题与致谢均译，参考文献按学术
        惯例保留英文著录与原文编号。
  \item \textbf{公式编号}：原文所有带编号公式一律以 \verb|\tag{...}| 锁定原文编号（如式 (3.21)），
        因此译本可与原文逐式对照；编号不随增删而改变，这是刻意的选择。
  \item \textbf{术语}：以中文学界约定俗成的译法为准（如 gauge theory 译“规范场论”、holonomy 译
        “和乐群”、anomaly 译“反常”），术语首次出现时附英文原词。
  \item \textbf{「本节重点」}（蓝色方框，标题后注“译者所加”）列出每节的核心物理图像、数学工具与
        结论；\textbf{「译注」}（橙色方框，全文连续编号）是学术审校批注，分“易错点、含糊点、物理图像、
        原文疑误”四类。两类内容均非原文，可分别在 \verb|main.tex| 中关闭。
  \item \textbf{原文笔误}：遇原文疑似笔误时一律照原文保留公式，另加\textbf{原文疑误}译注说明理由与
        建议，不擅自改动。
  \item \textbf{图表}：原文只有一张未编号的插图（第 31 页、式 (4.24) 与 (4.25) 之间的亏格 $g$ 黎曼面
        示意图）；译本按原文位置以图 1 插入，并在图注中注明原文页码与“原文未编号”。
  \item \textbf{编译}：需用 XeLaTeX 编译（\verb|xelatex main.tex| 两遍）。中文由 \verb|ctex| 宏包处理，
        默认字体集 \verb|fontset=fandol|；Windows 用户可改为 \verb|fontset=windows|。
\end{enumerate}

\clearpage
\tableofcontents
\clearpage
